\begin{explanationbox}

	\textbf{\textcolor{CExplanation}{一句话直觉}}

	两个等腰三角形共顶点且顶角相等时，将一个绕顶点任意旋转，则两非公共顶点的连线与对应顶点构成全等三角形，且这两条连线的夹角与顶角相等或互补。

	\medskip
	\textbf{\textcolor{CExplanation}{核心拆解}}
	\begin{description}[style=nextline, leftmargin=2.8em, labelsep=0.5em, itemsep=0.45em]
		\item[\textbf{要点一（全等条件）：}]
			$\triangle ABC$ 和 $\triangle ADE$ 均为等腰三角形（$AB=AC$, $AD=AE$），且顶角相等（$\angle BAC=\angle DAE=\alpha$）。这为旋转后的全等提供了 SAS 准备。
		\item[\textbf{要点二（旋转构造）：}]
			将 $\triangle ADE$ 绕点 $A$ 旋转任意角度，旋转后 $D,E$ 位置变化，但 $AD=AE$ 不变，$\angle DAE=\alpha$ 不变。连接 $BD$ 与 $CE$。
		\item[\textbf{要点三（全等结论）：}]
			由 SAS 得 $\triangle ABD\cong\triangle ACE$，故 $BD=CE$。进一步，$BD$ 与 $CE$ 的夹角等于 $\alpha$ 或 $180^\circ-\alpha$（取决于旋转方向）。
	\end{description}

	\medskip
	\textbf{\textcolor{CExplanation}{几何本质（旋转全等）}}

	手拉手的核心是旋转全等：两个等腰三角形在公共顶点处形成两组相等的边（$AB=AC$, $AD=AE$）和一组相等的角（$\angle BAD=\angle CAE$，借助 $\alpha$ 的加减得到）。旋转操作保证了对应边 $BD$ 和 $CE$ 的夹角与顶角关系固定。

	\medskip
	\textbf{\textcolor{CExplanation}{代数意义（坐标验证）}}

	在坐标系中，设 $A$ 为原点，$B,C,D,E$ 在单位圆上（或任意圆），利用旋转矩阵或复数乘法可证明 $BD=CE$ 和夹角关系，体现代数与几何的对应。

	\medskip
	\textbf{\textcolor{CExplanation}{考点价值（常见考法）}}
	\begin{itemize}[leftmargin=*, itemsep=0.5em, topsep=0.4em]
		\item \textbf{考法一（线段相等）：} 在复杂图形中识别手拉手模型，利用全等快速证明两线段相等，无需通过解三角形。
		\item \textbf{考法二（夹角计算）：} 已知顶角 $\alpha$ 和旋转角度，求 $BD$ 与 $CE$ 的夹角；或利用夹角关系反推旋转角度。
	\end{itemize}

	\medskip
	\textbf{\textcolor{CExplanation}{顿悟点}}

	理解“任意旋转”是关键：只要等腰和顶角相等保持不变，无论 $\triangle ADE$ 如何旋转，全等关系始终成立。这为动态几何问题提供了不变的等量关系。

	\medskip
	\textbf{\textcolor{CExplanation}{使用场景}}
	\begin{itemize}[leftmargin=*, itemsep=0.5em, topsep=0.4em]
		\item \textbf{场景一（图形识别）：} 当图中出现两个共顶点的等腰三角形（如正方形、等腰直角三角形、等边三角形）时，可考虑手拉手模型。
		\item \textbf{场景二（动态问题）：} 旋转、翻折等动态变换中，求某点运动路径长度或线段夹角，常需手拉手全等转化。
	\end{itemize}

\end{explanationbox}
